\section{Quasi-coherent sheaves on affines}
\label{section-quasi-coherent-affine}

\noindent
Recall that we have defined the abstract notion of a quasi-coherent
sheaf in Modules, Definition \ref{modules-definition-quasi-coherent}.
In this section we show that any quasi-coherent sheaf on an affine
scheme $\Spec(R)$ corresponds to the sheaf $\widetilde M$ associated to
an $R$-module $M$.

\begin{lemma}
\label{lemma-compare-constructions}
Let $(X, \mathcal{O}_X) = (\Spec(R), \mathcal{O}_{\Spec(R)})$
be an affine scheme. Let $M$ be an $R$-module. There exists a canonical
isomorphism between the sheaf $\widetilde M$ associated to the $R$-module
$M$ (Definition \ref{definition-structure-sheaf}) and the sheaf
$\mathcal{F}_M$ associated to the $R$-module $M$
(Modules, Definition \ref{modules-definition-sheaf-associated}).
This isomorphism is functorial in $M$. In particular,
the sheaves $\widetilde M$ are quasi-coherent. Moreover, they
are characterized by the following mapping property
$$
\Hom_{\mathcal{O}_X}(\widetilde M, \mathcal{F})
=
\Hom_R(M, \Gamma(X, \mathcal{F}))
$$
for any sheaf of $\mathcal{O}_X$-modules $\mathcal{F}$.
Here a map $\alpha : \widetilde M \to \mathcal{F}$ corresponds
to its effect on global sections.
\end{lemma}

\begin{proof}
By Modules, Lemma \ref{modules-lemma-construct-quasi-coherent-sheaves}
we have a morphism $\mathcal{F}_M \to \widetilde M$ corresponding
to the map $M \to \Gamma(X, \widetilde M) = M$. Let $x \in X$
correspond to the prime $\mathfrak p \subset R$.
The induced map on stalks are the maps
$\mathcal{O}_{X, x} \otimes_R M \to M_{\mathfrak p}$
which are isomorphisms because
$R_{\mathfrak p} \otimes_R M = M_{\mathfrak p}$.
Hence the map $\mathcal{F}_M \to \widetilde M$ is an isomorphism.
The mapping property follows from the mapping property of
the sheaves $\mathcal{F}_M$.
\end{proof}

\begin{lemma}
\label{lemma-widetilde-constructions}
Let $(X, \mathcal{O}_X) = (\Spec(R), \mathcal{O}_{\Spec(R)})$
be an affine scheme. There are canonical isomorphisms
\begin{enumerate}
\item
$
\widetilde{M \otimes_R N}
\cong
\widetilde M \otimes_{\mathcal{O}_X} \widetilde N
$,
see Modules, Section \ref{modules-section-tensor-product}.
\item
$
\widetilde{\text{T}^n(M)}
\cong
\text{T}^n(\widetilde M)
$,
$
\widetilde{\text{Sym}^n(M)}
\cong
\text{Sym}^n(\widetilde M)
$, and
$
\widetilde{\wedge^n(M)}
\cong
\wedge^n(\widetilde M)
$,
see
Modules, Section \ref{modules-section-symmetric-exterior}.
\item if $M$ is a finitely presented $R$-module, then
$
\SheafHom_{\mathcal{O}_X}(\widetilde M, \widetilde N)
\cong
\widetilde{\Hom_R(M,  N)}
$,
see
Modules, Section \ref{modules-section-internal-hom}.
\end{enumerate}
\end{lemma}

\begin{proof}[First proof]
Using Lemma \ref{lemma-compare-constructions} and
Modules, Lemma \ref{modules-lemma-construct-quasi-coherent-sheaves}
we see that the functor $M \mapsto \widetilde M$ can be viewed
as $\pi^*$ for a morphism $\pi$ of ringed spaces.
And pulling back modules commutes with tensor constructions by
Modules, Lemmas \ref{modules-lemma-tensor-product-pullback}
and \ref{modules-lemma-pullback-tensor-algebra}.
The morphism $\pi : (X, \mathcal{O}_X) \to (\{*\}, R)$ is
flat for example because the stalks of $\mathcal{O}_X$ are
localizations of $R$ (Lemma \ref{lemma-spec-sheaves})
and hence flat over $R$. Thus pullback by $\pi$ commutes
with internal hom if the first module is finitely presented by
Modules, Lemma \ref{modules-lemma-pullback-internal-hom}.
\end{proof}

\begin{proof}[Second proof]
Proof of (1). By Lemma \ref{lemma-compare-constructions} to give a map
$\widetilde{M \otimes_R N}$ into
$\widetilde M \otimes_{\mathcal{O}_X} \widetilde N$
we have to give a map on global sections
$M \otimes_R N \to
\Gamma(X, \widetilde M \otimes_{\mathcal{O}_X} \widetilde N)$
which exists by definition of the tensor product of sheaves
of modules. To see that this map is an isomorphism it
suffices to check that it is an isomorphism on stalks.
And this follows from the description of the stalks
of $\widetilde{M}$ (either in Lemma \ref{lemma-spec-sheaves} or in
Modules, Lemma \ref{modules-lemma-construct-quasi-coherent-sheaves}),
the fact that tensor product commutes with localization
(Algebra, Lemma \ref{algebra-lemma-tensor-product-localization}) and
Modules, Lemma \ref{modules-lemma-stalk-tensor-product}.

\medskip\noindent
Proof of (2). This is similar to the proof of (1), using
Algebra, Lemma \ref{algebra-lemma-tensor-algebra-localization} and
Modules, Lemma \ref{modules-lemma-stalk-tensor-algebra}.

\medskip\noindent
Proof of (3). Since the construction $M \mapsto \widetilde{M}$ is functorial
there is an $R$-linear map
$\Hom_R(M, N) \to \Hom_{\mathcal{O}_X}(\widetilde{M}, \widetilde{N})$.
The target of this map is the global sections of
$\SheafHom_{\mathcal{O}_X}(\widetilde M, \widetilde N)$.
Hence by Lemma \ref{lemma-compare-constructions} we obtain a map
of $\mathcal{O}_X$-modules $\widetilde{\Hom_R(M,  N)} \to
\SheafHom_{\mathcal{O}_X}(\widetilde M, \widetilde N)$.
We check that this is an isomorphism by comparing stalks.
If $M$ is finitely presented as an $R$-module
then $\widetilde M$ has a global finite presentation as an
$\mathcal{O}_X$-module. Hence we conclude using
Algebra, Lemma \ref{algebra-lemma-hom-from-finitely-presented} and
Modules, Lemma \ref{modules-lemma-stalk-internal-hom}.
\end{proof}

\begin{proof}[Third proof of part (1)]
For any $\mathcal{O}_X$-module $\mathcal{F}$
we have the following isomorphisms functorial in $M$, $N$, and $\mathcal{F}$
\begin{align*}
\Hom_{\mathcal{O}_X}(\widetilde{M} \otimes _{\mathcal{O} _X} \widetilde{N},
\mathcal{F})
& =
\Hom_{\mathcal{O}_X}(\widetilde{M},
\SheafHom_{\mathcal{O} _X} (\widetilde{N}, \mathcal{F})) \\
& =
\Hom_R(M, \Gamma(X,
\SheafHom_{\mathcal{O}_X}(\widetilde{N}, \mathcal{F}))) \\
& =
\Hom_R(M, \Hom_{\mathcal{O}_X}(\widetilde{N}, \mathcal{F})) \\
& =
\Hom_R(M, \Hom_R(N, \Gamma(X,\mathcal{F}))) \\
& =
\Hom_R(M \otimes_R N, \Gamma(X, \mathcal{F})) \\
& =
\Hom_{\mathcal{O}_X}(\widetilde{M \otimes_R N}, \mathcal{F})
\end{align*}
The first equality is Modules, Lemma \ref{modules-lemma-internal-hom}.
The second equality is the universal property of $\widetilde{M}$, see
Lemma \ref{lemma-compare-constructions}. The third equality holds by
definition of $\SheafHom$. The fourth equality is the universal property
of $\widetilde{N}$. Then fifth equality is
Algebra, Lemma \ref{algebra-lemma-hom-from-tensor-product}.
The final equality is the universal property of
$\widetilde{M \otimes_R N}$.
By the Yoneda lemma (Categories, Lemma \ref{categories-lemma-yoneda})
we obtain (1).
\end{proof}

\begin{lemma}
\label{lemma-widetilde-pullback}
Let
$(X, \mathcal{O}_X) = (\Spec(S), \mathcal{O}_{\Spec(S)})$,
$(Y, \mathcal{O}_Y) = (\Spec(R), \mathcal{O}_{\Spec(R)})$
be affine schemes.
Let $\psi : (X, \mathcal{O}_X) \to (Y, \mathcal{O}_Y)$ be a
morphism of affine schemes, corresponding to the ring map
$\psi^\sharp : R \to S$ (see Lemma \ref{lemma-category-affine-schemes}).
\begin{enumerate}
\item We have $\psi^* \widetilde M = \widetilde{S \otimes_R M}$
functorially in the $R$-module $M$.
\item We have $\psi_* \widetilde N = \widetilde{N_R}$ functorially
in the $S$-module $N$.
\end{enumerate}
\end{lemma}

\begin{proof}
The first assertion follows from the identification in
Lemma \ref{lemma-compare-constructions}
and the result of Modules, Lemma \ref{modules-lemma-restrict-quasi-coherent}.
The second assertion follows from the fact
that $\psi^{-1}(D(f)) = D(\psi^\sharp(f))$ and hence
$$
\psi_* \widetilde N(D(f)) = \widetilde N(D(\psi^\sharp(f))) =
N_{\psi^\sharp(f)} = (N_R)_f = \widetilde{N_R}(D(f))
$$
as desired.
\end{proof}

\noindent
Lemma \ref{lemma-widetilde-pullback} above says in particular
that if you restrict
the sheaf $\widetilde M$ to a standard affine open subspace
$D(f)$, then you get $\widetilde{M_f}$. We will use this from
now on without further mention.

\begin{lemma}
\label{lemma-quasi-coherent-affine}
Let $(X, \mathcal{O}_X) = (\Spec(R), \mathcal{O}_{\Spec(R)})$
be an affine scheme. Let $\mathcal{F}$ be a
quasi-coherent $\mathcal{O}_X$-module. Then
$\mathcal{F}$ is isomorphic to the sheaf associated to
the $R$-module $\Gamma(X, \mathcal{F})$.
\end{lemma}

\begin{proof}
Let $\mathcal{F}$ be a quasi-coherent $\mathcal{O}_X$-module.
Since every standard open $D(f)$ is quasi-compact we see that
$X$ is locally quasi-compact, i.e., every point has a fundamental
system of quasi-compact neighbourhoods, see Topology,
Definition \ref{topology-definition-locally-quasi-compact}.
Hence by Modules, Lemma \ref{modules-lemma-quasi-coherent-module}
for every prime $\mathfrak p \subset R$ corresponding to $x \in X$
there exists an open neighbourhood $x \in U \subset X$ such that
$\mathcal{F}|_U$ is isomorphic to the quasi-coherent
sheaf associated to some $\mathcal{O}_X(U)$-module $M$.
In other words, we get an open covering by $U$'s with this property.
By Lemma \ref{lemma-standard-open} for example we can refine this
covering to  a standard open covering.
Thus we get a covering $\Spec(R) = \bigcup D(f_i)$
and $R_{f_i}$-modules $M_i$ and isomorphisms
$\varphi_i : \mathcal{F}|_{D(f_i)} \to \mathcal{F}_{M_i}$
for some $R_{f_i}$-module $M_i$. On the overlaps
we get isomorphisms
$$
\xymatrix{
\mathcal{F}_{M_i}|_{D(f_if_j)}
\ar[rr]^{\varphi_i^{-1}|_{D(f_if_j)}}
& &
\mathcal{F}|_{D(f_if_j)}
\ar[rr]^{\varphi_j|_{D(f_if_j)}}
& &
\mathcal{F}_{M_j}|_{D(f_if_j)}.
}
$$
Let us denote these $\psi_{ij}$. It is clear that
we have the cocycle condition
$$
\psi_{jk}|_{D(f_if_jf_k)}
\circ
\psi_{ij}|_{D(f_if_jf_k)}
=
\psi_{ik}|_{D(f_if_jf_k)}
$$
on triple overlaps.

\medskip\noindent
Recall that each of the open subspaces $D(f_i)$, $D(f_if_j)$,
$D(f_if_jf_k)$ is an affine scheme. Hence the sheaves $\mathcal{F}_{M_i}$
are isomorphic to the sheaves $\widetilde M_i$ by Lemma
\ref{lemma-compare-constructions} above. In particular we see that
$\mathcal{F}_{M_i}(D(f_if_j)) = (M_i)_{f_j}$, etc.
Also by Lemma \ref{lemma-compare-constructions} above we see
that $\psi_{ij}$ corresponds to a unique $R_{f_if_j}$-module isomorphism
$$
\psi_{ij} : (M_i)_{f_j} \longrightarrow (M_j)_{f_i}
$$
namely, the effect of $\psi_{ij}$ on sections over $D(f_if_j)$.
Moreover these then satisfy the cocycle condition that
$$
\xymatrix{
(M_i)_{f_jf_k}
\ar[rd]_{\psi_{ij}}
\ar[rr]^{\psi_{ik}}
& &
(M_k)_{f_if_j} \\
&
(M_j)_{f_if_k} \ar[ru]_{\psi_{jk}}
}
$$
commutes (for any triple $i, j, k$).

\medskip\noindent
Now Algebra, Lemma \ref{algebra-lemma-glue-modules}
shows that there exist an $R$-module $M$ such that
$M_i = M_{f_i}$ compatible with the morphisms $\psi_{ij}$.
Consider $\mathcal{F}_M = \widetilde M$. At this point it is
a formality to show that $\widetilde M$ is isomorphic to
the quasi-coherent sheaf $\mathcal{F}$ we started out with.
Namely, the sheaves $\mathcal{F}$ and $\widetilde M$ give
rise to isomorphic sets of glueing data of sheaves of $\mathcal{O}_X$-modules
with respect to the covering $X = \bigcup D(f_i)$, see
Sheaves, Section \ref{sheaves-section-glueing-sheaves}
and in particular Lemma \ref{sheaves-lemma-mapping-property-glue}.
Explicitly, in the current situation, this boils down to
the following argument: Let us construct an $R$-module map
$$
M \longrightarrow \Gamma(X, \mathcal{F}).
$$
Namely, given $m \in M$ we get $m_i = m/1 \in M_{f_i} = M_i$
by construction of $M$. By construction of $M_i$ this corresponds
to a section $s_i \in \mathcal{F}(U_i)$. (Namely, $\varphi^{-1}_i(m_i)$.)
We claim that $s_i|_{D(f_if_j)} = s_j|_{D(f_if_j)}$. This is
true because, by construction of $M$, we have $\psi_{ij}(m_i) = m_j$,
and by the construction of the $\psi_{ij}$. By the sheaf condition of
$\mathcal{F}$ this collection of sections gives rise to a unique
section $s$ of $\mathcal{F}$ over $X$. We leave it to the reader
to show that $m \mapsto s$ is a $R$-module map.
By Lemma \ref{lemma-compare-constructions} we obtain an associated
$\mathcal{O}_X$-module map
$$
\widetilde M \longrightarrow \mathcal{F}.
$$
By construction this map reduces to the isomorphisms
$\varphi_i^{-1}$ on each $D(f_i)$ and hence is an isomorphism.
\end{proof}

\begin{lemma}
\label{lemma-equivalence-quasi-coherent}
Let $(X, \mathcal{O}_X) = (\Spec(R), \mathcal{O}_{\Spec(R)})$
be an affine scheme.
The functors $M \mapsto \widetilde M$ and
$\mathcal{F} \mapsto \Gamma(X, \mathcal{F})$ define quasi-inverse
equivalences of categories
$$
\xymatrix{
\QCoh(\mathcal{O}_X)
\ar@<1ex>[r]
&
\text{Mod}_R
\ar@<1ex>[l]
}
$$
between the category of quasi-coherent $\mathcal{O}_X$-modules
and the category of $R$-modules.
\end{lemma}

\begin{proof}
See Lemmas \ref{lemma-compare-constructions}
and \ref{lemma-quasi-coherent-affine} above.
\end{proof}

\noindent
From now on we will not distinguish between quasi-coherent
sheaves on affine schemes and sheaves of the form $\widetilde M$.

\begin{lemma}
\label{lemma-kernel-cokernel-quasi-coherent}
Let $X = \Spec(R)$ be an affine scheme.
Kernels and cokernels of maps of quasi-coherent
$\mathcal{O}_X$-modules are quasi-coherent.
\end{lemma}

\begin{proof}
This follows from the exactness of the functor $\widetilde{\ }$
since by Lemma \ref{lemma-compare-constructions} we know that any map
$\psi : \widetilde{M} \to \widetilde{N}$ comes from
an $R$-module map $\varphi : M \to N$. (So we have
$\Ker(\psi) = \widetilde{\Ker(\varphi)}$ and
$\Coker(\psi) = \widetilde{\Coker(\varphi)}$.)
\end{proof}

\begin{lemma}
\label{lemma-colimit-quasi-coherent}
Let $X = \Spec(R)$ be an affine scheme.
The direct sum of an arbitrary collection of quasi-coherent sheaves
on $X$ is quasi-coherent. The same holds for colimits.
\end{lemma}

\begin{proof}
Suppose $\mathcal{F}_i$, $i \in I$ is a collection of quasi-coherent
sheaves on $X$. By Lemma \ref{lemma-equivalence-quasi-coherent}
above we can write $\mathcal{F}_i = \widetilde{M_i}$ for some $R$-module
$M_i$. Set $M = \bigoplus M_i$. Consider the sheaf $\widetilde{M}$.
For each standard open $D(f)$ we have
$$
\widetilde{M}(D(f)) = M_f =
\left(\bigoplus M_i\right)_f =
\bigoplus M_{i, f}.
$$
Hence we see that the quasi-coherent $\mathcal{O}_X$-module
$\widetilde{M}$ is the direct sum of the sheaves $\mathcal{F}_i$.
A similar argument works for general colimits.
\end{proof}

\begin{lemma}
\label{lemma-extension-quasi-coherent}
Let $(X, \mathcal{O}_X) = (\Spec(R), \mathcal{O}_{\Spec(R)})$
be an affine scheme. Suppose that
$$
0 \to
\mathcal{F}_1 \to
\mathcal{F}_2 \to
\mathcal{F}_3 \to
0
$$
is a short exact sequence of sheaves of $\mathcal{O}_X$-modules.
If two out of three are quasi-coherent then so is the third.
\end{lemma}

\begin{proof}
This is clear in case both $\mathcal{F}_1$ and $\mathcal{F}_2$ are
quasi-coherent because the functor $M \mapsto \widetilde M$
is exact, see Lemma \ref{lemma-spec-sheaves}.
Similarly in case both $\mathcal{F}_2$ and $\mathcal{F}_3$ are
quasi-coherent. Now, suppose that $\mathcal{F}_1 = \widetilde M_1$ and
$\mathcal{F}_3 = \widetilde M_3$ are quasi-coherent.
Set $M_2 = \Gamma(X, \mathcal{F}_2)$. We claim it suffices to show that
the sequence
$$
0 \to M_1 \to M_2 \to M_3 \to 0
$$
is exact. Namely, if this is the case, then (by using the mapping
property of Lemma \ref{lemma-compare-constructions}) we get a commutative
diagram
$$
\xymatrix{
0 \ar[r] &
\widetilde M_1 \ar[r] \ar[d] &
\widetilde M_2 \ar[r] \ar[d] &
\widetilde M_3 \ar[r] \ar[d] &
0 \\
0 \ar[r] &
\mathcal{F}_1 \ar[r] &
\mathcal{F}_2 \ar[r] &
\mathcal{F}_3 \ar[r] &
0
}
$$
and we win by the snake lemma.

\medskip\noindent
The ``correct'' argument here would be to show first
that $H^1(X, \mathcal{F}) = 0$ for any quasi-coherent sheaf $\mathcal{F}$.
This is actually not all that hard, but it is perhaps better to postpone
this till later. Instead we use a small trick.

\medskip\noindent
Pick $m \in M_3 = \Gamma(X, \mathcal{F}_3)$.
Consider the following set
$$
I = \{ f \in R \mid \text{the element }fm\text{ comes from }M_2\}.
$$
Clearly this is an ideal. It suffices to show $1 \in I$.
Hence it suffices to show that for any prime $\mathfrak p$
there exists an $f \in I$, $f \not\in \mathfrak p$.
Let $x \in X$ be the point corresponding to $\mathfrak p$.
Because surjectivity can be checked on stalks
there exists an open neighbourhood $U$ of $x$ such that
$m|_U$ comes from a local section $s \in \mathcal{F}_2(U)$.
In fact we may assume that $U = D(f)$ is a standard open,
i.e., $f \in R$, $f \not \in \mathfrak p$. We will show
that for some $N \gg 0$ we have $f^N \in I$, which
will finish the proof.

\medskip\noindent
Take any point $z \in V(f)$, say corresponding to the
prime $\mathfrak q \subset R$. We can also find a $g \in R$,
$g \not \in \mathfrak q$ such that $m|_{D(g)}$ lifts
to some $s' \in \mathcal{F}_2(D(g))$.
Consider the difference $s|_{D(fg)} - s'|_{D(fg)}$.
This is an element $m'$ of $\mathcal{F}_1(D(fg)) = (M_1)_{fg}$.
For some integer $n = n(z)$ the element $f^n m'$ comes
from some $m'_1 \in (M_1)_g$. We see that
$f^n s$ extends to a section $\sigma$ of $\mathcal{F}_2$ on $D(f) \cup D(g)$
because it agrees with the restriction of
$f^n s' + m'_1$ on $D(f) \cap D(g) = D(fg)$.
Moreover, $\sigma$ maps to the restriction of $f^n m$
to $D(f) \cup D(g)$.

\medskip\noindent
Since $V(f)$ is quasi-compact, there exists a finite list
of elements $g_1, \ldots, g_m \in R$ such that
$V(f) \subset \bigcup D(g_j)$, an integer $n > 0$ and sections
$\sigma_j \in \mathcal{F}_2(D(f) \cup D(g_j))$ such that
$\sigma_j|_{D(f)} = f^n s$ and $\sigma_j$ maps to the section
$f^nm|_{D(f) \cup D(g_j)}$ of $\mathcal{F}_3$.
Consider the differences
$$
\sigma_j|_{D(f) \cup D(g_jg_k)}
-
\sigma_k|_{D(f) \cup D(g_jg_k)}.
$$
These correspond to sections of $\mathcal{F}_1$
over $D(f) \cup D(g_jg_k)$ which are zero
on $D(f)$. In particular their images in
$\mathcal{F}_1(D(g_jg_k)) = (M_1)_{g_jg_k}$
are zero in $(M_1)_{g_jg_kf}$.
Thus some high power of $f$ kills each and every one of these.
In other words, the elements $f^N \sigma_j$, for some $N \gg 0$
satisfy the glueing condition of the sheaf property and
give rise to a section $\sigma $ of $\mathcal{F}_2$
over $\bigcup (D(f) \cup D(g_j)) = X$ as desired.
\end{proof}







